\chapter{Introduction générale}

\section{Introduction}
Dans le cadre de la transformation digitale des organisations, les départements de développement jouent un rôle essentiel dans la conception d'outils internes adaptés aux besoins métier. Au sein de Royal Air Maroc, la Digital Factory constitue un département dédié au développement et à la réalisation de solutions numériques. Elle accueille également des stagiaires qui participent à des projets applicatifs encadrés par des responsables techniques.

Ce chapitre présente le contexte général du stage, l'organisme d'accueil, la problématique traitée, les objectifs du projet et l'organisation globale du rapport. Il permet de situer clairement la différence entre la Digital Factory, qui est le département d'accueil au sein de Royal Air Maroc, et la solution développée, qui est une plateforme de gestion des projets de stage.

\section{Contexte}
La Digital Factory de Royal Air Maroc accueille plusieurs stagiaires et porte des initiatives numériques variées. Le suivi des stages nécessite une coordination continue entre les responsables, les encadrants et les stagiaires : définition des projets, affectation des personnes, planification des tâches, réception des livrables et suivi de l'avancement.

Lorsque ces informations sont dispersées entre différents outils de communication et espaces de stockage, la visibilité sur les projets devient plus difficile. Le présent projet vise ainsi à fournir un point d'entrée unique pour organiser et suivre les projets de stage.

\section{Problématique}
Le suivi des projets de stage repose sur plusieurs intervenants et sur des informations qui évoluent tout au long du stage. Sans outil centralisé, l'affectation des stagiaires, le suivi des tâches, le partage des livrables et les retours des encadrants peuvent être dispersés entre différents supports. Cette situation limite la visibilité sur l'avancement global et complique la coordination entre les membres de l'équipe.

Comment concevoir une plateforme web légère et sécurisée permettant de centraliser la gestion des projets de stage, de fluidifier les échanges entre encadrants et stagiaires, et de donner une vision claire de l'état d'avancement des travaux ?

\section{Objectifs}
L'objectif principal est de développer une plateforme web de gestion des projets de stage pour la Digital Factory de Royal Air Maroc. Les objectifs spécifiques sont les suivants :
\begin{itemize}
  \item gérer les comptes des encadrants et des stagiaires ;
  \item créer des projets et attribuer les stagiaires ainsi que les encadrants ;
  \item planifier, prioriser et suivre les tâches ;
  \item centraliser les soumissions, pièces jointes et liens externes ;
  \item faciliter les commentaires et le retour des encadrants ;
  \item assister le SuperAdmin dans la préparation des projets grâce à l'IA ;
  \item intégrer la qualité et la sécurité dans le cycle de livraison par une démarche DevSecOps.
\end{itemize}

\section{Organisme d'accueil}
Royal Air Maroc (RAM) est la compagnie aérienne nationale marocaine. Dans un environnement où les services numériques jouent un rôle central, l'innovation et l'amélioration continue des outils internes constituent des leviers importants pour accompagner les équipes et les projets. La Digital Factory est le département d'accueil du stage au sein de RAM. Elle contribue à la conception et à la réalisation de solutions numériques, notamment à travers des projets de développement applicatif. Dans ce cadre, l'encadrement de plusieurs stagiaires et de leurs projets appelle des mécanismes simples de coordination, de traçabilité et de partage des livrables. La plateforme développée pendant ce stage répond à ce besoin en proposant une solution interne qui relie les rôles de gouvernance, d'encadrement et de réalisation, tout en conservant l'historique des activités liées aux projets et aux tâches.

\section{Organisation du rapport}
Ce premier chapitre présente le contexte, l'organisme d'accueil et les objectifs du projet. Le deuxième chapitre détaille l'analyse fonctionnelle, la conception et la planification. Le troisième chapitre présente la réalisation technique, la chaîne DevSecOps et les interfaces de la plateforme. Le rapport se termine par une conclusion générale et des perspectives d'évolution.

\section{Conclusion}
Ce chapitre a permis de situer le projet dans son contexte professionnel et organisationnel. La Digital Factory a été présentée comme le département d'accueil au sein de Royal Air Maroc, tandis que la solution développée a été définie comme une plateforme web dédiée à la gestion des projets de stage. Les chapitres suivants présentent l'analyse, la conception, la planification puis la réalisation technique de cette plateforme.
